\section{Introduction}
\label{sec:introduction}

\PARstart{D}{ay-ahead} load forecasts are a central input of power system operation. Transmission system operators (TSOs) use them to schedule generation, set reserve requirements, and plan transmission constraints, so forecast errors translate directly into operating costs and security margins~\cite{hong2016tutorial,hong2020outlook}. Two developments are changing this task. First, the rapid diffusion of behind-the-meter photovoltaic (PV) generation makes the demand seen by the operator increasingly weather dependent and continuously reshapes the daily load profile, a persistent source of concept drift~\cite{gama2014drift}. Second, open data policies now give researchers access to the same measurements that operators use, which makes it possible to compare openly reproducible models with the forecasts actually used in operation, something rarely done in the load forecasting literature~\cite{hong2020outlook,petropoulos2022forecasting}.

These developments raise three practical questions that the literature does not answer jointly. (i) How accurate are open, data-driven day-ahead models, built only with information available at the operator's decision time, relative to the official forecast of a TSO? (ii) When load is drifting, how much history should be used to fit a model and how often should it be refitted? In electricity price forecasting, the choice and combination of calibration windows have been shown to matter~\cite{marcjasz2018windows,hubicka2019windows}, but for load this question has received little attention outside the abrupt COVID-19 shock~\cite{obst2021covid}. (iii) Can prediction intervals remain calibrated as the distribution of errors changes, for instance through adaptive conformal methods~\cite{gibbs2021aci,zaffran2022aci}?

We address these questions using a system that offers an unusually clear case of both developments. The Brazilian National Interconnected System (SIN) serves about 80~GW of average demand over a continental territory divided into four subsystems with distinct climates and consumption patterns. After the legal framework for distributed micro- and mini-generation (MMGD) was consolidated~\cite{lei14300}, behind-the-meter generation, mostly rooftop PV, grew from about 5\% of the midday load in 2021 to roughly one third in 2026, according to the operator's own estimates (Section~\ref{sec:data}). The National System Operator (ONS) publishes the hourly load of each subsystem since 2000, its MMGD estimates, and the official day-ahead load forecast used in the daily operation programming~\cite{ons_opendata,ons_apicarga}. The case therefore allows a long, fully reproducible evaluation against an operational reference, under a drift of a magnitude that other systems with fast PV adoption are likely to face.

The main contributions of this paper are as follows:
\begin{itemize}
\item We release an open and fully reproducible benchmark for day-ahead hourly load forecasting of the four SIN subsystems, built exclusively from public operator data (2012--2026), with an operationally realistic information set (forecasts issued at noon of the previous day, horizons of 14--37~h), a test period of almost three years never used for model design, and statistical significance testing.
\item We provide what is, to our knowledge, the first systematic comparison between open forecasting models and the official day-ahead forecast of a TSO under rapid behind-the-meter PV growth, including how the relative performance varies by hour of the day, by subsystem, and over time.
\item We quantify how the calibration window length (1 year to the full history) and the refitting frequency (daily to yearly) affect accuracy under drift, for linear and tree-based models, and show that a simple level normalisation of the target is decisive for tree ensembles.
\item We compare quantile gradient boosting, rolling empirical (split-conformal) intervals, and adaptive conformal inference with delayed feedback for day-ahead probabilistic load forecasting, and assess their calibration over time.
\item We analyse errors on special days (national holidays, Carnival, bridge days, and the year-end period), which are a major source of operational error in Brazil and many other countries~\cite{ziel2018holidays}.
\end{itemize}

The remainder of the paper is organised as follows. Section~\ref{sec:related} reviews related work. Section~\ref{sec:data} describes the data and Section~\ref{sec:method} the forecasting framework and evaluation protocol. Section~\ref{sec:results} presents the results, Section~\ref{sec:discussion} discusses their implications and limitations, and Section~\ref{sec:conclusion} concludes.
